\section{A logarithmic unfinished-block bound under conditioning}
\label{sec:conditioned-clock}
The stationary law of an induced block is length-biased. We use this law
for the initial block and the exact visit intensity for subsequent blocks.
These two ingredients avoid applying a launch-law union bound to a
stationary initial state.

Write $\overline\tau_R^*=\bar\tau_R/c_*$ and represent the flow as the
suspension over $(Y_R^*,F_R^*,\nu_R^*)$ with roof $\tau_R^*$. Let $t_j$
be its two-sided visit times and $x_j$ their outgoing section states.
The minimum time between consecutive visits is at least $3/50$.

\begin{lemma}[Marked visit intensity]
\label{lem:marked-visit-intensity}
For every nonnegative measurable $f$ on $Y_R^*$ and every $t\ge0$,
\begin{equation}\label{eq:marked-visit-intensity}
 \mathbb E_{\mu_R}\sum_{t_j\in(0,t]} f(x_j)
       =\frac{t}{\overline\tau_R^*}\int f\dd\nu_R^*.
\end{equation}
\end{lemma}
\begin{proof}
First take bounded $f$. Stationarity makes the locally finite measure
$A\mapsto\mathbb E_{\mu_R}\sum_{t_j\in A}f(x_j)$ translation-invariant
on $\R$, hence equal to $i_f$ times Lebesgue measure. Local finiteness
follows from the positive lower roof bound. Choose $0<\delta<3/50$.
At a stationary time, being in the first $\delta$ units of an induced
roof is the same event as a visit in the preceding interval of length
$\delta$, apart from null endpoints. There is at most one such visit,
and its mark is the current base state. Integration in the suspension
therefore gives
\[
 \delta i_f=\frac1{\overline\tau_R^*}
       \int_{Y_R^*}\int_0^\delta f(x)\dd s\dd\nu_R^*(x)
       =\frac{\delta}{\overline\tau_R^*}\int f\dd\nu_R^*.
\]
This determines $i_f$ and proves the formula for bounded $f$.
Monotone convergence extends it to arbitrary nonnegative marks, including
the case of an infinite integral. No renewal independence is involved.
\end{proof}

Let $\mathcal M_R(t)$ be the maximum of $Q_R$ over all complete induced
blocks whose time intervals meet $[0,t]$, including the block containing
time zero. Endpoint conventions at a visit are immaterial for $\mu_R$.
The marked size of any unfinished piece is bounded by a fixed multiple
of $1+\mathcal M_R(t)$.

\begin{theorem}[Uniform exponential maximum and conditioned residual]
\label{thm:conditioned-maximal-block}
There are constants $a_2,C_2>0$ such that, for every $R\in I$, $t\ge0$
and $q\ge0$,
\begin{equation}\label{eq:exponential-maximal-block}
 \mu_R\{\mathcal M_R(t)>q\}\le C_2(1+t)e^{-a_2q}.
\end{equation}
Consequently, for any measurable event $B_{R,t}$ of probability at least
$d_t>0$,
\begin{equation}\label{eq:conditioned-maximal-block}
 \mu_R\{\mathcal M_R(t)>q\mid B_{R,t}\}
       \le\min\{1,C_2(1+t)e^{-a_2q}/d_t\}.
\end{equation}
In particular, if $d_t\ge d(1+t)^{-A}$ with $d>0$ and $A\ge0$, then
$\mathcal M_R(t)=O_P(\log(2+t))$ under these conditional laws, uniformly
in $R$ and in the events. The maximum marked unfinished return during
$[0,t]$ is then $o_P(\sqrt t)$ under the same laws. More generally the
last assertion holds if $\log((1+t)/d_t)=o(\sqrt t)$.
\end{theorem}
\begin{proof}
The initial containing block has law
$\tau_R^*\dd\nu_R^*/\overline\tau_R^*$. If the maximum exceeds $q$,
either that block has $Q_R>q$ or a visit in $(0,t]$ starts a block with
$Q_R>q$. The union bound and Lemma~\ref{lem:marked-visit-intensity} give
\begin{equation}\label{eq:exact-maximal-tail-budget}
 \mu_R\{\mathcal M_R(t)>q\}
 \le\frac1{\overline\tau_R^*}
       \left(\int\tau_R^*\one_{\{Q_R>q\}}\dd\nu_R^*
                       +t\,\nu_R^*(Q_R>q)\right).
\end{equation}
The denominator is bounded below on $I$. Since $\tau_R^*\le Q_R$,
the exponential moment in \eqref{eq:exponential-return} bounds the first
integral by $C e^{-a_1q/2}$ and the probability by $C e^{-a_1q}$.
For example, $Q\le (2/a_1)e^{a_1Q/2}$ and
$e^{a_1Q/2}\one_{Q>q}\le e^{a_1Q}e^{-a_1q/2}$ suffice for the
first bound. This proves \eqref{eq:exponential-maximal-block}.

Dividing the unconditional exceptional probability by $\mu_R(B_{R,t})$
proves \eqref{eq:conditioned-maximal-block}. With a polynomial lower
bound on $d_t$, setting $q=K\log(2+t)$ and taking
$a_2K>A+1$ makes the right side tend to zero. With the more general
hypothesis, set $q=\epsilon\sqrt t$ for any $\epsilon>0$; the logarithm
of the displayed upper bound tends to minus infinity. The deterministic
bound of each marked piece by $C(1+Q_R)$ proves the assertions for all
unfinished pieces at once. The event may depend on the whole orbit and
on an estimator from that orbit; no independence is assumed.
\end{proof}

For a stationary starting point $(x,s)$ put
$V_R(v)=\max\{j:T_{j,R}(x)\le s+v\}$. Define the two path records
\[
 Z_R(v)=(V_R(v),C_R(v),W_R(v)),\qquad
 Z_R^\circ(v)=(V_R(v),N_{V_R(v),R}(x),K_{V_R(v),R}(x)).
\]
The second record counts full blocks from the beginning of the initial
containing roof, not from the observation time. It is used solely as an
exact intermediate record for comparison.

\begin{corollary}[A conditional path-comparison budget]
\label{cor:conditional-path-budget}
There is $C_3$ such that, pathwise outside the collision singularities,
\begin{equation}\label{eq:deterministic-prefix-bound}
 \sup_{0\le v\le t}|Z_R(v)-Z_R^\circ(v)|
                         \le C_3(1+\mathcal M_R(t)).
\end{equation}
Let $\widehat Z_R$ and $\widehat Z_R^\circ$ be the corresponding paths
on $[0,1]$, divided by $\sqrt t$ and with any identical deterministic
centering. For $|F|\le1$ Lipschitz with constant $L_F$ for the uniform
path metric and $\mu_R(B_{R,t})\ge d_t>0$,
\begin{align}\label{eq:conditional-path-budget}
 \left|\mathbb E_{\mu_R}[F(\widehat Z_R)\mid B_{R,t}]
       -\mathbb E_{\mu_R}[F(\widehat Z_R^\circ)\mid B_{R,t}]\right|
 &\le\frac{C_3L_F(1+q)}{\sqrt t}
        +\frac{2C_2(1+t)e^{-a_2q}}{d_t}.
\end{align}
For $t\ge2$ choose
\begin{equation}\label{eq:conditional-log-budget}
 q=a_2^{-1}\log\left(\frac{\max\{1,2C_2\}(1+t)\sqrt t}{d_t}\right).
\end{equation}
The second term is at most $t^{-1/2}$; for polynomially bounded
conditioning probabilities the total error is
$O((1+L_F)\log t/\sqrt t)$.
\end{corollary}
\begin{proof}
The collision counts differ by the initial prefix and the final
incomplete block. Each is bounded by the full block's collision count.
The same holds for the lattice displacement after multiplying by the
uniform bound on a single-flight lattice step and adding a fixed
fundamental-cell boundary constant. The visit counts agree exactly by
construction. These estimates are simultaneous for all $v\in[0,t]$,
and prove \eqref{eq:deterministic-prefix-bound}. On
$\{\mathcal M_R(t)\le q\}$ the Lipschitz bound gives the first term in
\eqref{eq:conditional-path-budget}. On its complement the difference of
two values of $F$ is at most two. Apply
\eqref{eq:conditioned-maximal-block}. The stated choice of $q$ and
substitution prove the final assertions.
\end{proof}

\begin{remark}[Which conditional interface has been closed]
\label{rem:conditional-interface-scope}
The square-root unfinished-block requirement is now a theorem for the
actual mechanical family, even after any conditioning of polynomially
small positive probability and over the whole observation interval. A
functional central limit theorem for completed records is still a
separate premise in Proposition~\ref{prop:clock}. For a local limit, the
conditioning event on the two sides of \eqref{eq:conditional-path-budget}
is the same event. Replacing an exact lattice observation by the
completed-block observation changes that event and needs a further
boundary estimate. If two events $B,B'$ have probabilities at least $d$,
then for $|f|\le1$ direct subtraction of the ratios gives
\[
 |\mathbb E[f\mid B]-\mathbb E[f\mid B']|
             \le2\mu_R(B\mathbin\triangle B')/d.
\]
Thus an $O(\log t)$ error in an integer record is not discarded when
conditioning on one lattice value. The weighted raw-density estimates
of Proposition~\ref{prop:conditioning} remain the requisite local input.
A zero-probability singleton in the continuous coordinate is not an event
to which \eqref{eq:conditioned-maximal-block} applies.
\end{remark}
